% problem_id: S001-lemma-structure-decomposition
% source_id: S001 (Stacks Project)
% source_file: pione.tex
% statement_locator: pione.tex lines 3174--3210
% proof_locator: pione.tex lines 3212--3342
% env: lemma
% id_note: label 在本来源内唯一
% pairing: tight (verified)
% status: complete (statement + author proof verbatim + reason + locators)
%
\begin{lemma}
\label{lemma-structure-decomposition}
Let $A$ be a discrete valuation ring with fraction field $K$.
Let $L/K$ be a (possibly infinite) Galois extension.
Let $B$ be the integral closure of $A$ in $L$.
Let $\mathfrak m$ be a maximal ideal of $B$.
Let $G = \text{Gal}(L/K)$,
$D = \{\sigma \in G \mid \sigma(\mathfrak m) = \mathfrak m\}$, and
$I = \{\sigma \in D \mid \sigma \bmod \mathfrak m =
\text{id}_{\kappa(\mathfrak m)}\}$.
The decomposition group $D$ fits into a canonical exact sequence
$$
1 \to I \to D \to \text{Aut}(\kappa(\mathfrak m)/\kappa_A) \to 1
$$
The inertia group $I$ fits into a canonical exact sequence
$$
1 \to P \to I \to I_t \to 1
$$
such that
\begin{enumerate}
\item $P$ is a normal subgroup of $D$,
\item $P$ is a pro-$p$-group if the characteristic of
$\kappa_A$ is $p > 1$ and $P = \{1\}$ if the characteristic of $\kappa_A$
is zero,
\item there is a multiplicatively directed $S \subset \mathbf{N}$
such that $\kappa(\mathfrak m)$ contains a primitive $n$th root of unity
for each $n \in S$ (elements of $S$ are prime to $p$),
\item there exists a canonical surjective map
$$
\theta_{can} : I \to \lim_{n \in S} \mu_n(\kappa(\mathfrak m))
$$
whose kernel is $P$, which satisfies
$\theta_{can}(\tau \sigma \tau^{-1}) = \tau(\theta_{can}(\sigma))$
for $\tau \in D$, $\sigma \in I$, and which induces an isomorphism
$I_t \to \lim_{n \in S} \mu_n(\kappa(\mathfrak m))$.
\end{enumerate}
\end{lemma}

\begin{proof}
This is mostly a reformulation of the results on finite Galois extensions
proved in More on Algebra, Section \ref{more-algebra-section-ramification}.
The surjectivity of the map $D \to \text{Aut}(\kappa(\mathfrak m)/\kappa)$ is
More on Algebra, Lemma \ref{more-algebra-lemma-one-orbit-geometric-galois}.
This gives the first exact sequence.

\medskip\noindent
To construct the second short exact sequence let $\Lambda$ be the set
of finite Galois subextensions, i.e., $\lambda \in \Lambda$ corresponds
to $L/L_\lambda/K$. Set $G_\lambda = \text{Gal}(L_\lambda/K)$.
Recall that $G_\lambda$ is an inverse system of finite groups with surjective
transition maps and that $G = \lim_{\lambda \in \Lambda} G_\lambda$, see
Fields, Lemma \ref{fields-lemma-infinite-galois-limit}.
We let $B_\lambda$ be the integral closure of $A$ in $L_\lambda$.
Then we set $\mathfrak m_\lambda = \mathfrak m \cap B_\lambda$
and we denote $P_\lambda, I_\lambda, D_\lambda$ the
wild inertia, inertia, and decomposition group of
$\mathfrak m_\lambda$, see More on Algebra, Lemma
\ref{more-algebra-lemma-galois-inertia}.
For $\lambda \geq \lambda'$ the restriction defines
a commutative diagram
$$
\xymatrix{
P_\lambda \ar[d] \ar[r] &
I_\lambda \ar[d] \ar[r] &
D_\lambda \ar[d] \ar[r] &
G_\lambda \ar[d] \\
P_{\lambda'} \ar[r] &
I_{\lambda'} \ar[r] &
D_{\lambda'} \ar[r] &
G_{\lambda'}
}
$$
with surjective vertical maps, see
More on Algebra, Lemma \ref{more-algebra-lemma-compare-inertia}.

\medskip\noindent
From the definitions it follows immediately
that $I = \lim I_\lambda$ and $D = \lim D_\lambda$
under the isomorphism $G = \lim G_\lambda$ above.
Since $L = \colim L_\lambda$ we have $B = \colim B_\lambda$
and $\kappa(\mathfrak m) = \colim \kappa(\mathfrak m_\lambda)$.
Since the transition maps of the system $D_\lambda$
are compatible with the maps
$D_\lambda \to \text{Aut}(\kappa(\mathfrak m_\lambda)/\kappa)$
(see More on Algebra, Lemma \ref{more-algebra-lemma-compare-inertia})
we see that the map $D \to \text{Aut}(\kappa(\mathfrak m)/\kappa)$
is the limit of the maps
$D_\lambda \to \text{Aut}(\kappa(\mathfrak m_\lambda)/\kappa)$.

\medskip\noindent
There exist canonical maps
$$
\theta_{\lambda, can} :
I_\lambda
\longrightarrow
\mu_{n_\lambda}(\kappa(\mathfrak m_\lambda))
$$
where $n_\lambda = |I_\lambda|/|P_\lambda|$, where
$\mu_{n_\lambda}(\kappa(\mathfrak m_\lambda))$ has
order $n_\lambda$, such that
$\theta_{\lambda, can}(\tau \sigma \tau^{-1}) =
\tau(\theta_{\lambda, can}(\sigma))$ for
$\tau \in D_\lambda$ and $\sigma \in I_\lambda$, and such that
we get commutative diagrams
$$
\xymatrix{
I_\lambda \ar[r]_-{\theta_{\lambda, can}} \ar[d] &
\mu_{n_\lambda}(\kappa(\mathfrak m_\lambda))
\ar[d]^{(-)^{n_\lambda/n_{\lambda'}}} \\
I_{\lambda'} \ar[r]^-{\theta_{\lambda', can}} &
\mu_{n_{\lambda'}}(\kappa(\mathfrak m_{\lambda'}))
}
$$
see
More on Algebra, Remark \ref{more-algebra-remark-canonical-inertia-character}.

\medskip\noindent
Let $S \subset \mathbf{N}$ be the collection of integers $n_\lambda$.
Since $\Lambda$ is directed, we see that $S$ is multiplicatively directed.
By the displayed commutative diagrams above we can take the limits of
the maps $\theta_{\lambda, can}$ to obtain
$$
\theta_{can} : I \to \lim_{n \in S} \mu_n(\kappa(\mathfrak m)).
$$
This map is continuous (small detail omitted). Since the transition maps
of the system of $I_\lambda$ are surjective
and $\Lambda$ is directed, the projections $I \to I_\lambda$
are surjective. For every $\lambda$ the diagram
$$
\xymatrix{
I \ar[d] \ar[r]_-{\theta_{can}} &
\lim_{n \in S} \mu_n(\kappa(\mathfrak m)) \ar[d] \\
I_{\lambda} \ar[r]^-{\theta_{\lambda, can}} &
\mu_{n_\lambda}(\kappa(\mathfrak m_\lambda))
}
$$
commutes. Hence the image of $\theta_{can}$ surjects onto the finite group
$\mu_{n_\lambda}(\kappa(\mathfrak m)) =
\mu_{n_\lambda}(\kappa(\mathfrak m_\lambda))$ of order $n_\lambda$
(see above). It follows that the image of $\theta_{can}$ is dense.
On the other hand $\theta_{can}$ is continuous and the
source is a profinite group. Hence $\theta_{can}$ is surjective
by a topological argument.

\medskip\noindent
The property $\theta_{can}(\tau \sigma \tau^{-1}) = \tau(\theta_{can}(\sigma))$
for $\tau \in D$, $\sigma \in I$ follows from the corresponding properties
of the maps $\theta_{\lambda, can}$ and the compatibility of the map
$D \to \text{Aut}(\kappa(\mathfrak m))$ with the maps
$D_\lambda \to \text{Aut}(\kappa(\mathfrak m_\lambda))$.
Setting $P = \Ker(\theta_{can})$ this implies
that $P$ is a normal subgroup of $D$. Setting $I_t = I/P$
we obtain the isomorphism $I_t \to \lim_{n \in S} \mu_n(\kappa(\mathfrak m))$
from the surjectivity of $\theta_{can}$.

\medskip\noindent
To finish the proof we show that $P = \lim P_\lambda$ which proves
that $P$ is a pro-$p$-group. Recall that the tame inertia group
$I_{\lambda, t} = I_\lambda/P_\lambda$ has order $n_\lambda$.
Since the transition maps $P_\lambda \to P_{\lambda'}$ are surjective
and $\Lambda$ is directed, we obtain a short exact sequence
$$
1 \to \lim P_\lambda \to I \to \lim I_{\lambda, t} \to 1
$$
(details omitted). Since for each $\lambda$ the map $\theta_{\lambda, can}$
induces an isomorphism
$I_{\lambda, t} \cong \mu_{n_\lambda}(\kappa(\mathfrak m))$
the desired result follows.
\end{proof}
